\section{Mechanical consequences of the threshold coefficients}
\label{sec:consequences}
\subsection{A multi-roof effect not determined by the one-roof law}
\begin{corollary}[The nonrenewal onset exponent]\label{cor:nonrenewal}
Fix $R\in(0,1/2)$. Let an independent-roof stationary suspension have
the same section roof marginal as the Lorentz gas. Its intensity is
therefore the same $\lambda_R$. For each fixed $j\ge1$, its maximal
count onset satisfies
\begin{equation}\label{eq:renewal}
 \Pr^{\rm ind}(N_{jg_R+\varepsilon}=j+1)
 =\frac{\lambda_R}{(j+1)!}
       \left(\frac3{\sqrt{g_R}}\right)^j
       \varepsilon^{j+1}\bigl(1+O_j(\varepsilon)\bigr).
\end{equation}
For the actual Lorentz gas the exponent is two for every $j$, with
coefficient \eqref{eq:main-onset}. In particular the two laws have
different leading powers for every $j\ge2$, despite their identical
one-roof marginal and intensity.
\end{corollary}
\begin{proof}
The case $j=1$ of the action calculation gives
$\int(g_R+\varepsilon-\tau_R)_+\dd\nu
=3\varepsilon^2/(2\sqrt{g_R})+O(\varepsilon^3)$.
The smooth threshold density proved above permits differentiation.
Thus the excess roof $u=\tau_R-g_R$ has density
$3/\sqrt{g_R}+O(u)$ near $u=0$.
For $j$ independent excess roofs, integrate their product density on
$u_i\ge0$, $\sum_i u_i<\varepsilon$ with weight
$\varepsilon-\sum_i u_i$. The leading integral is
$\varepsilon^{j+1}/(j+1)!$, as follows by integrating successively
on this simplex. The error is $O_j(\varepsilon)$ relative to it.
The stationary identity \eqref{eq:tails-new}, with
$\varepsilon<g_R$, supplies the factor $\lambda_R$ and excludes the
next count. This proves \eqref{eq:renewal}. The actual law is
Theorem~\ref{thm:unweighted}.
\end{proof}

For example the first genuinely multi-roof count has
\begin{equation}\label{eq:three-hit}
 \Pr_R(N_{2g_R+\varepsilon}=3)
 =\frac{3\lambda_R R}{4(1-R)\sqrt{g_R}}
       \varepsilon^2\bigl(1+O(\varepsilon)\bigr).
\end{equation}
The independent-roof substitute instead has coefficient
$3\lambda_R/(2g_R)$ multiplying $\varepsilon^3$.
The stationary identity is the same in both systems; the difference
comes from the specular boundary-value geometry of $S_2\tau_R$.

\subsection{Recovery of periodic-orbit instability}
\begin{corollary}[Instability in count probabilities]\label{cor:multiplier}
The unstable multiplier of the normal period-two return map is
$e^{2\chi_R}$. The measured threshold coefficients satisfy
\begin{equation}\label{eq:ratio}
 \lim_{j\to\infty}\frac{C_{j+1}(R)}{C_j(R)}=e^{-\chi_R},\qquad
 \frac{C_1(R)}{C_2(R)}=2\cosh\chi_R.
\end{equation}
Thus they determine that multiplier. The unweighted coefficient series
has the exact representation
\begin{equation}\label{eq:lambert}
 \sum_{j\ge1}C_j(R)z^j
 =\frac{3\lambda_R}{R}
       \sum_{m\ge0}\frac{z e^{-(2m+1)\chi_R}}
                         {1-z e^{-(2m+1)\chi_R}},\qquad |z|<e^{\chi_R}.
\end{equation}
It extends meromorphically to the complex plane with simple poles
at $z=e^{(2m+1)\chi_R}$, $m\ge0$.
\end{corollary}
\begin{proof}
The Jacobi recurrence in the proof of Lemma~\ref{lem:bridge} has
transfer matrix $\begin{psmallmatrix}2c_R&-1\\1&0\end{psmallmatrix}$
and eigenvalues $e^{\chi_R},e^{-\chi_R}$. The local change between
successive transverse positions and a collision position--momentum
pair is invertible since the one-flight twist is $1/g_R$.
After two flights the coordinate frame returns to the original
orientation, so the physical period-two derivative is conjugate to
the square of this transfer matrix. This proves the multiplier claim.
The ratios follow from the explicit $C_j$.

Expand $\csch(j\chi)=2\sum_{m\ge0}e^{-(2m+1)j\chi}$ and sum the
absolutely convergent geometric series in $j$. On every compact set
avoiding the displayed poles, the tail in $m$ converges uniformly,
since its terms are $O(|z|e^{-(2m+1)\chi})$. This gives the
meromorphic continuation. At each indicated point only one summand
has a pole and its residue is nonzero, so there is no cancellation.
\end{proof}

Series \eqref{eq:lambert} organizes threshold coefficients belonging to
different observation windows. It is not the ordinary pressure or the
generating function of counts in one fixed window. The recovery of
instability here uses equilibrium tail amplitudes; it does not claim
that the relation between periodic action Hessians and monodromy is new.

\subsection{The retained regular-window response}
The complete companion \cite{TC} proves its fixed-window theorem,
including the all-order moving-level formula, without using the present
threshold asymptotic. We record the terminal-level witness requested
for that calculation. At $T=11/100$, $R\in[9/20,47/100]$, choose
$\alpha=0$ facing $(1,0)$ and
\begin{equation}\label{eq:terminal-witness}
 \cos\vartheta=\frac{1-2R+T^2}{2T(1-R)}.
\end{equation}
The right side lies strictly between zero and one. The point
$(R,0)+T(\cos\vartheta,\sin\vartheta)$ lies on the target circle,
since substitution gives squared distance $R^2$ from its center.
It is an incoming hit because
\[
 (1-R)\cos\vartheta-T=\frac{1-2R-T^2}{2T}>0.
\]
The outgoing sign and the no-obstruction estimate in the companion
make it a genuine first hit. The terminal level is therefore nonempty.
Its uniform gradient bound and $\partial_R\tau<0$ show directly that
the level integral in the companion's second-derivative formula is
strictly positive. This is compatible with the onset singularity:
the fixed window remains strictly above $g_R$ on its specified interval,
whereas \eqref{eq:main-onset} crosses the critical level itself.
